\chapter{Modélisation}

\section{Projet}

Un projet contient l'ensemble des données nécessaires à l'analyse :

\begin{itemize}
  \item géométrie ;
  \item matériaux ;
  \item sections ;
  \item appuis ;
  \item charges ;
  \item combinaisons ;
  \item paramètres de calcul.
\end{itemize}

Les projets sont sauvegardés dans un fichier SQLite avec l'extension
\fichier{.db}.

\section{Nœuds}

Les nœuds définissent la géométrie de base du modèle. Chaque nœud possède :

\begin{itemize}
  \item un numéro unique ;
  \item une coordonnée X ;
  \item une coordonnée Y ;
  \item une coordonnée Z ;
  \item des conditions aux limites éventuelles.
\end{itemize}

Pour ajouter un nœud par coordonnées, utiliser le menu ou le bouton
\bouton{Ajouter un nœud}. Le numéro choisi doit être unique.

\begin{infobox}
Si les coordonnées saisies correspondent déjà à un nœud existant, le logiciel
refuse la création afin d'éviter deux nœuds superposés.
\end{infobox}

\section{Barres}

Une barre relie deux nœuds. Elle nécessite au minimum :

\begin{itemize}
  \item un nœud initial ;
  \item un nœud final ;
  \item une section ;
  \item un matériau associé à cette section.
\end{itemize}

Les barres sont utilisées pour modéliser poutres, poteaux et éléments linéaires.

\section{Surfaces}

Les surfaces permettent de modéliser des plaques ou dalles. Une surface peut
être triangulaire ou quadrangulaire selon la formulation choisie.

\begin{warningbox}
Les surfaces sont calculées uniquement avec les fonctionnalités actuellement
prises en charge par le moteur sélectionné. Si le moteur courant ne supporte pas
les plaques, le logiciel désactive les actions correspondantes.
\end{warningbox}

\section{Grille}

La grille facilite la saisie et le dessin en définissant des lignes selon les
axes X, Y et Z. Elle peut être utilisée pour des modèles 3D complets ou pour des
plans de travail 2D.

\todoDoc{Ajouter une capture du dialogue de grille et un exemple de grille de portique.}
